\section{Constraints Can Clarify Strategy}

\subsection{Hard, Negotiable, and Assumed Boundaries}

Founders are often encouraged to think without limits. Imagination benefits from freedom, but operating decisions do not. A business is built within a set of constraints: time, financial obligations, regulations, employment, geography, capability, infrastructure, and personal commitments.

A useful distinction separates 3 classes of boundaries:

\begin{itemize}[leftmargin=11pt,itemsep=1.5pt,topsep=1.5pt,parsep=0pt]
  \item \textbf{Hard constraints:} Violation would be illegal, unethical, financially destructive, or incompatible with obligations you must honor. Examples: using employer code, selling accredited credentials without legal authority, or risking emergency reserves.
  \item \textbf{Negotiable constraints:} Expensive or inconvenient, but changeable with effort. Examples: lack of recording gear, no physical office, or small ad budget.
  \item \textbf{Assumed constraints:} Things you have never tested. Examples: \textit{``No corporate buyer will pay without a famous logo''} or \textit{``I must write 10,000 lines of software before offering a consulting diagnostic.''}
\end{itemize}

\textbf{Distinguish hard constraints from assumed constraints: never gamble essential obligations on unverified assumptions.}

\subsection{The Constraint Register}

\begin{tcolorbox}[
  enhanced,
  colback=white,
  colframe=BayesBorder!80!white,
  boxrule=0.4pt,
  arc=2pt,
  left=4pt,right=4pt,top=4pt,bottom=4pt,
  before skip=3pt,after skip=5pt
]
\scriptsize
\setlength{\tabcolsep}{2.5pt}
\renewcommand{\arraystretch}{1.18}
\begin{tabularx}{\linewidth}{@{} >{\bfseries\color{BayesInk}}p{1.8cm} >{\centering\arraybackslash}p{1.1cm} >{\color{BayesInk}}X >{\color{BayesSecondary}}p{2.6cm} @{}}
\toprule
\textbf{Constraint} & \textbf{Type} & \textbf{Consequence} & \textbf{Investigation / Path} \\
\midrule
No employer IP or data & Hard & Many product shortcuts are off-limits & Build de-novo with open or client-licensed data. \\
\midrule
Evenings-only live availability & Hard (at present) & Cannot service 9--5 enterprise helpdesks & Run cohort sessions or asynchronous tooling. \\
\midrule
Zero institutional alumni brand & Current reality & Buyers face higher perceived risk & Offer unbundled pilots, verifiable student code. \\
\midrule
``Clients demand local offices'' & Assumed belief & Risk of premature commercial lease & Test remote delivery and client on-sites first. \\
\bottomrule
\end{tabularx}
\end{tcolorbox}

In some cases a constraint is a source of product design. Limited instructor hours may encourage better asynchronous diagnostics and a narrower curriculum. Limited capital may encourage pre-selling a small cohort rather than building an elaborate platform.

\subsection{The Reversibility Filter}

Before taking an action, ask how easily you can undo it.

A 1-month pilot with modest cash outlay is highly reversible. Signing a multi-year commercial lease, resigning before validating demand, hiring a large permanent team, or taking personal debt against speculative future tuition is dangerous and nearly impossible to undo.

\textbf{Uncertain evidence should buy reversible experiments, not irreversible commitments.} Exceptions exist when scarce opportunities require decisive action, but they must be argued explicitly rather than hidden behind blind optimism.
